
量子交换的母模型不应从单模开始。先假设环境谱已经由上一章的一般 Green/正常模分析选出 $M$ 个需要显式保留的正交玻色模式，但电子仍保留完整正能连续谱。第 $\nu$ 个模式的中心频率为 $\omega_\nu$，湮灭算符为 $\hat a_\nu$。在相互作用绘景中，最一般的线性单量子交换仍是连续核
\begin{equation}
\boxed{
\hat H_{{\rm Q},I}^{(M)}(t)
=
\hbar
\sum_{\nu=1}^{M}
\sum_{ss'}
\int\dd^3p\,\dd^3p'\,
\left[
\mathcal G_{\nu;s's}(\mathbf p',\mathbf p;t)
|\mathbf p',s'\rangle\langle\mathbf p,s|
\,\hat a_\nu
+\mathrm{H.c.}
\right].}
\label{eq:multimode_qpinem_continuous_kernel_dp101}
\end{equation}
$\mathcal G_{\nu;s's}$ 由 QED 电流--模式矩阵元或相应材料量子化给出；这一层还没有电子“格点”，也没有为每个光学模式预先指定唯一动量步长。

很多纳米光子学几何中，第 $\nu$ 个模式只在一个窄的纵向 Fourier 区间内与准直电子强耦合。若实际电子工作窗内的主导转移为 $\hbar\kappa_\nu$，并且自旋翻转及其它动量转移的累计传播修正很小，可以在通道子空间上写
\[
\hat{\mathcal G}_{\nu,I}(t)
=
\ii\upsilon_\nu(t)\hat b_\nu^\dagger
+\hat R_{\nu,\kappa}(t),
\qquad
\hat b_\nu^\dagger|p\rangle
=
|p+\hbar\kappa_\nu\rangle .
\]
这里 $\hat R_{\nu,\kappa}$ 收集非主导 Fourier 分量、逐边矩阵元变化和被舍弃自旋通道。若 $\Pi_{\rm in}$ 投影到实际入射工作窗，以只保留主导平移项的传播子记为 $\hat U_{\rm tr}$，则一个直接的有限时间充分条件是
\begin{equation}
\boxed{
\epsilon_{\kappa}^{(M)}
\equiv
\frac1\hbar
\int_{t_i}^{t_f}\dd t\,
\left\|
\hbar\sum_\nu
\left[
\hat R_{\nu,\kappa}(t)\hat a_\nu
+\hat R_{\nu,\kappa}^\dagger(t)\hat a_\nu^\dagger
\right]
\hat U_{\rm tr}(t,t_i)\Pi_{\rm in}
\right\|
\ll1 .}
\label{eq:multimode_momentum_step_projection_error_dp101}
\end{equation}
只有通过这一类通道投影后，才得到常用的多模平移 Hamiltonian
\begin{equation}
\boxed{
 \hat H_{\rm Q}^{(M)}(t)
 =\ii\hbar\sum_{\nu=1}^{M}
 \left[
 \upsilon_\nu(t)\hat b_\nu^\dagger\hat a_\nu
 -\upsilon_\nu^*(t)\hat b_\nu\hat a_\nu^\dagger
 \right].}
\label{eq:multimode_qpinem_hamiltonian_dp47}
\end{equation}
因此 $\upsilon_\nu$ 由连续耦合核在选定电子通道上的投影唯一确定。

重复交换以后，整数向量
\[
\boldsymbol\ell=(\ell_1,\ldots,\ell_M)\in\mathbb Z^M
\]
映射到真实电子动量
\begin{equation*}
\mathbf p_{\boldsymbol\ell}
=
\mathbf p_0
+\hbar\sum_\nu\ell_\nu\boldsymbol\kappa_\nu .
\end{equation*}
把 $|\boldsymbol\ell\rangle_e$ 作为正交的多维格点标签，还要求整数映射在实际占据窗口 $\mathcal W_\ell$ 内保持一一对应。若存在不同整数向量满足
\[
\sum_\nu
(\ell_\nu-\ell_\nu')\boldsymbol\kappa_\nu=0,
\]
它们代表同一个电子动量末态，必须先在振幅层面相干相加，而不能作为不同格点分别计数。对可唯一标记的通道，真实电子能量为
\begin{equation}
 \mathcal E_{\boldsymbol\ell}
 =E\!\left(
 \mathbf p_0+\hbar\sum_\nu
 \ell_\nu\boldsymbol\kappa_\nu
 \right),
\end{equation}
并相对于理想多频能量梯定义
\begin{equation}
\boxed{
 \delta_{\boldsymbol\ell}
 =
 \frac{
 \mathcal E_{\boldsymbol\ell}-E_0
 -\sum_\nu\ell_\nu\hbar\omega_\nu
 }{\hbar}.}
\label{eq:multimode_recoil_detuning_dp47}
\end{equation}

对初始 Fock 扇区
$\mathbf n=(n_1,\ldots,n_M)$，取守恒基底
\[
 |\boldsymbol\ell;\mathbf n\rangle
 =
 |\boldsymbol\ell\rangle_e
 \otimes
 \bigotimes_\nu
 |n_\nu-\ell_\nu\rangle_\nu ,
\]
其中每个光子占据数都必须非负。把联合态展开为
\[
|\Psi_{\mathbf n}(t)\rangle
=
\sum_{\boldsymbol\ell}
c_{\boldsymbol\ell}^{(\mathbf n)}(t)
|\boldsymbol\ell;\mathbf n\rangle ,
\]
逐项使用玻色梯矩阵元后得到
\begin{equation}
\boxed{
\begin{aligned}
 \ii\dot c_{\boldsymbol\ell}^{(\mathbf n)}
 =&\,\delta_{\boldsymbol\ell}
 c_{\boldsymbol\ell}^{(\mathbf n)}\\
 &+\ii\sum_\nu \upsilon_\nu(t)
 \sqrt{n_\nu-\ell_\nu+1}\,
 c_{\boldsymbol\ell-\mathbf e_\nu}^{(\mathbf n)}\\
 &-\ii\sum_\nu \upsilon_\nu^*(t)
 \sqrt{n_\nu-\ell_\nu}\,
 c_{\boldsymbol\ell+\mathbf e_\nu}^{(\mathbf n)} .
\end{aligned}}
\label{eq:multimode_qpinem_lattice_dp47}
\end{equation}
这才是有限组显式模式下的多模 QPINEM 数值一般动力学方程：真实电子色散、多个 Fock 下界、有限脉冲和模式间共享的电子反冲全部保留。单模模型只需令 $M=1$，但若进一步要得到双边位移闭式，还必须额外满足逐边失谐与逐边耦合近似均匀。

对多模相干态
$|\{\alpha_\nu\}\rangle$，经典场极限应逐模取
\[
|\alpha_\nu|\to\infty,
\qquad
\gqmode{\nu}\to0,
\qquad
\beta_\nu
\equiv
\alpha_\nu\gqmode{\nu}
\quad\text{保持有限}.
\]
若在相应工作窗中各电子平移算符可用完整连续谱上的平移表示，它们彼此对易，电子得到
\begin{equation}
\boxed{
 S_{\rm cl}^{(M)}
 =
 \exp\!\left[
 \sum_\nu
 \left(
 \beta_\nu\hat b_\nu^\dagger
 -\beta_\nu^*\hat b_\nu
 \right)
 \right].}
\label{eq:multicolor_classical_scattering_dp47}
\end{equation}
若各动量步长进一步满足
$\kappa_\nu=m_\nu\kappa_0$，不同颜色的整数路径可能到达同一个最终电子格点。把每个颜色的 Jacobi--Anger 展开相乘，并对所有满足
$\ell=\sum_\nu m_\nu r_\nu$
的路径作相干求和，得到
\begin{equation}
\boxed{
 A_\ell
 =\sum_{\{r_\nu\}}
 \delta_{\ell,\sum_\nu m_\nu r_\nu}
 \prod_\nu
 J_{r_\nu}(2|\beta_\nu|)
 \ee^{\ii r_\nu\varphi_\nu},
 \qquad
 \varphi_\nu=\arg\beta_\nu .}
\label{eq:multicolor_pinem_convolution_dp47}
\end{equation}
因此多色 PINEM 的相位敏感性来自“不同交换路径汇聚到同一电子末态”后的振幅干涉；若路径对应不同可分辨电子动量，则应保留多指标末态，而不能先把概率卷积到一条能量格上。
